\section{总结与展望}

\subsection{工作总结}

本文围绕视觉退化条件下 EEG 辅助视觉语言模型的语义识别与生成式描述问题, 在 Thought2Text/EEGCVPR 小样本数据集上完成了系统的研究. 本文的主要工作总结如下:

\textbf{(1) A2 temporal-spectral-spatial EEG 语义融合模型}是本文最强的定量模型. A2 通过三个独立分支并行提取 EEG 的时间、频谱和空间通道信息, 采用 gated residual fusion 在 CLIP 语义空间中融合图像和 EEG 嵌入. 在 3 种子 × 6 退化条件的 18 组对照评估中, real EEG 在全部条件下稳定超过 vision-only、shuffled EEG 和 random EEG, 证明了真实配对 EEG 在退化视觉条件下具有可机器利用的语义增强能力.

\textbf{(2) VTF token-level EEG-enhanced VLM}成功将 EEG 信息从 pooled embedding 层推进到 token 层. VTF 通过 visual-to-EEG cross-attention 将 EEG tokens 与 CLIP ViT visual tokens 融合, 生成 EEG-enhanced vision tokens. 该结构为后续将 EEG 信息输入生成式 VLM 提供了直接可行的通道路径.

\textbf{(3) Route5 Qwen2-VL generative EVLM 原型}将 EEG-enhanced vision tokens 输入 Qwen2-VL-2B-Instruct 生成式模型, 通过 LoRA ($r=8$) 参数高效微调, 在 strong degradation 条件下 real EEG 的 caption class-hit 为 37.36\% (best-of-N reranking) 和 38.26\% (LoRA r8), 比 vision-only 提高约 9 个百分点, 比 shuffled/random EEG 提高约 8--11 个百分点. Valid caption rate 在 best-of-N 配置下达到 97.75\%. 该原型验证了从小样本 EEG-image 数据构建生成式 EEG-enhanced VLM 的可行性.

\subsection{创新点}

\begin{enumerate}
  \item \textbf{构建了视觉退化条件下 EEG-VLM 语义辅助对照任务}: 系统定义了六种退化条件、五种输入模式和多维度评价指标, 为验证 EEG 语义增强作用提供了严格的实验框架;
  \item \textbf{设计了 A2 temporal-spectral-spatial EEG encoder}: 首次在 EEG 视觉语义解码中同时显式建模时间、频谱和空间通道三个维度的特征, 并通过 gated residual fusion 实现输入自适应的融合;
  \item \textbf{实现了 EEG-enhanced CLIP ViT visual tokens}: 通过 VTF cross-attention 将 EEG 信息在 token 层面融入视觉表示, 使 EEG 信号可以进入原生 VLM 的 token 处理流程;
  \item \textbf{构建了小样本 EEG 条件下的 Qwen2-VL 生成式 VLM 原型}: 完成了从 EEG 信号到自然语言 caption 的完整通路, 并进行了 LoRA rank、caption target 和 reranking 的系统消融.
\end{enumerate}

\subsection{不足与展望}

本文存在以下局限, 为未来研究方向提供了参考:

\textbf{数据集规模小}: Thought2Text 仅包含 40 类约 12K trial, 远小于典型 VLM 的预训练数据. 更大规模的 EEG-image-text 数据集（如 EIT-1M 的全量版本或合成增强数据）可能带来显著的性能提升.

\textbf{Caption supervision 弱}: 本文使用 BLIP 伪 caption 作为训练目标, 其质量远不如人工撰写的高质量 caption. 未来的工作可以使用更大更强的 captioning model（如 BLIP-2 或 LLaVA）生成更准确的伪 caption, 或引入部分人工标注的自然语言描述.

\textbf{EEG 个体差异未充分建模}: 本文所有训练和评估均在 6 个被试的混合数据上进行, 没有进行 subject-invariant 建模. 将 subject adaptation 模块（如 subject-specific adapter）加入 EEG encoder 可能显著提升跨被试的泛化能力.

\textbf{token fusion 的复杂度有限}: VTF 仅使用一次 cross-attention, 信息交互深度不足. 更复杂的 fusion 结构（如 multi-layer Q-Former 或 Perceiver IO）可能提升 EEG-enhanced vision tokens 的语义质量.

\textbf{生成式模型仍依赖 semantic prompt}: 当前 Route5 的生成质量和 class-hit 仍高度依赖输入中的 semantic top-k 提示信息. 完全依赖 EEG-enhanced vision tokens（不附加 prompt）进行自由生成是该方向的最终目标, 但当前尚未达到.

未来方向包括: (1) 在更大规模的 EEG-image-text 数据集上进行预训练和迁移; (2) 引入 subject adaptation 降低个体差异的影响; (3) 使用更强的 Q-Former/Perceiver bridge 替代简单的 cross-attention; (4) 探索更稳定的生成式训练目标, 减少生成模型对文本提示的依赖.
